\section{Conclusion}
\label{sec:conclusion}

Multi-turn agent RL requires credit assignment across long interaction trajectories, yet current frameworks default to step-independent windowing that decomposes trajectories into independently optimized windows. This default introduces a structural credit assignment pathology that no hyperparameter tuning can resolve.

We diagnose a dual mechanism underlying this failure: policy information loss at the inference layer, which imposes a constant performance ceiling, and advantage sign misattribution at the training layer, which causes cumulative degradation. Formal analysis shows that at a baseline success rate of $p = 12.5\%$ with group size $K = 8$, approximately 80.9\% of gradient updates for locally correct windows point in the wrong direction. Our window-size ablation reveals a credit assignment cliff: as $h$ decreases, performance does not degrade smoothly but collapses in two phases, with action validity dropping from 89.3\% to 43.1\%. Episode-level GRPO, which treats the full trajectory as a single optimization sample, eliminates both mechanisms and achieves 28.1\% validation success where step-independent training collapses to 0\%.

The central insight is that sign errors are more destructive than magnitude imprecision. Gradients that point in the wrong direction cause active divergence, while gradients with imprecise magnitude merely slow convergence. This distinction applies to any agent RL method that combines step-level windowing, sparse binary reward, and group-relative advantage normalization.

Several directions extend this work. First, selective verification trains the agent to learn when self-verification is warranted during execution. Second, cross-environment validation on WebShop and ScienceWorld would test whether empirical severity matches our theoretical predictions across different trajectory lengths and success rates. Third, scaling experiments with 7B and 14B models will reveal how sign misattribution interacts with model capacity. Finally, $\lambda$-annealing from episode-level ($\lambda = 1$) to finer-grained credit ($\lambda < 1$) offers a principled interpolation between direction correctness and credit granularity as the policy improves.
